\chapter{Wprowadzenie}

\section{Wstęp}
Na przestrzeni ostatnich kilku lat można zaobserwować wzrost w zainteresowaniu różnego rodzaju technikami wspomagającymi zdrowie psychiczne jak i fizyczne \cite{ristovski2023importance}. Powodem tego jest między innymi stale rosnące tempo życia oraz coraz większe wymagania zawodowe i edukacyjne. Popularną metodą pomagającą radzeniu sobie z daną presją mogą być sesje ćwiczeń oddechowych przeprowadzane za pomocą filmów instruktażowych lub aplikacji mobilnych. 

W przypadku aplikacji mobilnych pożądane jest również możliwość personalizacji oraz monitorowanie treningu, aby użytkownik mógł dostosować odpowiednio poziom trudności czy otrzymać informację zwrotną jak dana sesja przebiegła.

\section{Cele pracy}
Celem pracy jest opracowanie narzędzia wspomagającego realizację treningów oddechowych w postaci aplikacji mobilnej, która umożliwi użytkownikowi definiowanie indywidualnych wzorców oddechowych oraz zapewni monitorowanie poszczególnych faz oddechu w czasie trwania sesji treningowej. 

W ramach badań przeprowadzona zostanie analiza istniejących metod monitorowania oddechu, zarówno wykorzystujących dedykowane czujniki, jak i rozwiązania oparte na powszechnie dostępnych urządzeniach mobilnych. Na podstawie przeprowadzonego przeglądu literatury oraz istniejących rozwiązań zostanie zaprojektowana i zaimplementowana aplikacja mobilna umożliwiająca realizację treningów oddechowych z możliwością personalizacji wzorców oddechowych. Dodatkowo do monitorowania wzorców oddechu skonstruowany zostanie model sztucznej inteligencji wykorzystujący uczenie głębokie. Opracowane rozwiązanie zostanie następnie poddane weryfikacji pod kątem poprawności działania oraz przydatności w kontekście wspomagania treningów oddechowych.

Wyniki pracy mogą przyczynić się do zwiększenia dostępności narzędzi wspierających praktykę ćwiczeń oddechowych, a także umożliwić ich łatwiejszą integrację z codziennym życiem użytkowników.